\documentclass[12pt, a4paper, twocolumn]{article}

% --- PAQUETES NECESARIOS ---
\usepackage[utf8]{inputenc}
\usepackage[spanish]{babel}
\usepackage[margin=2.5cm]{geometry} % Márgenes exactos de 2,5 cm
\usepackage{times} % Fuente Times New Roman
\usepackage{titlesec} % Para cambiar el tamaño de los títulos
\usepackage{abstract} % Para personalizar el abstract

% --- CONFIGURACIÓN DE FORMATO SEGÚN EL PDF ---
\setlength{\columnsep}{1cm} % Espacio entre columnas de 1 cm
\pagestyle{empty} % Sin números de página ni encabezados


% Formato de secciones (Times New Roman, 12, negrita)
\titleformat{\section}{\normalfont\fontsize{12}{15}\bfseries}{\thesection.}{1em}{}

% --- TÍTULO Y AUTORES ---
% (Times New Roman, 16, negrita)
\title{\vspace{-2cm}\fontsize{16}{20}\bfseries Optimización de la Inclinación y Orientación de Paneles Solares Mediante Algoritmos Genéticos}

% Autores (Times New Roman, 14, negrita) e Institución (12, negrita, cursiva)
\author{
    \fontsize{14}{17}\bfseries Apellido/s, Nombre/s \\
    \fontsize{12}{15}\bfseries\itshape Universidad Tecnológica Nacional, Facultad Regional Buenos Aires
}
\date{} % Sin fecha para que no aparezca abajo


\begin{document}

\begin{titlepage}

\centering

{\fontsize{16}{20}\selectfont\bfseries
Optimización de la Inclinación y Orientación de Paneles Solares Mediante Algoritmos Genéticos\par}

\vspace{1.5cm}

\begin{minipage}{0.9\textwidth}
\centering
\begin{tabular}{p{0.45\textwidth} p{0.45\textwidth}}

\centering\textbf{Arach, Mateo Simón} &
\centering\textbf{Borsato Gimenez, Milton Rubén}\tabularnewline[0.5cm]

\centering Legajo: 51394 &
\centering Legajo: 51391\tabularnewline

\centering Correo: matearach@gmail.com &
\centering Correo: borsatomilton@gmail.com\tabularnewline[0.8cm]

\centering\textbf{Dayer, José Ignacio} &
\centering\textbf{Marchese, Valentin David}\tabularnewline[0.5cm]

\centering Legajo: 51508 &
\centering Legajo: 51745\tabularnewline

\centering Correo: nachodayer@gmail.com &
\centering Correo: marchese52002@gmail.com\tabularnewline

\end{tabular}
\end{minipage}

\vspace{1.5cm}

{\fontsize{12}{15}\selectfont\bfseries\itshape
Universidad Tecnológica Nacional\\
Facultad Regional Rosario
}

\vfill

\end{titlepage}

%==========================
% A partir de acá, dos columnas
%==========================


\twocolumn

\newpage

% --- CUERPO DEL DOCUMENTO ---
% Todo lo que sigue está en Times New Roman 12

\noindent\textbf{Palabras Clave}

Algoritmos genéticos, energía solar fotovoltaica, optimización, inclinación y orientación de paneles.

\section*{Introducción}
La energía solar fotovoltaica se ha consolidado como una de las principales alternativas para la transición hacia matrices energéticas sostenibles, debido a su bajo impacto ambiental, su costo decreciente de instalación y su capacidad de despliegue tanto en instalaciones residenciales como industriales. Sin embargo, la eficiencia real de un sistema fotovoltaico no depende únicamente de la calidad tecnológica de los paneles, sino también de factores geométricos y ambientales que condicionan la cantidad de radiación solar efectivamente captada. En particular, la inclinación y la orientación de los paneles son variables determinantes del rendimiento energético, y su cálculo mediante fórmulas empíricas o aproximaciones estáticas suele ignorar variaciones estacionales, condiciones climáticas locales y sombras del entorno. Este trabajo tiene como objetivo determinar si un algoritmo genético permite optimizar dichos parámetros para maximizar la generación energética de una instalación fotovoltaica.

\section*{Denominación del futuro proyecto de investigación}
Optimización de la inclinación y orientación de paneles solares fotovoltaicos mediante algoritmos genéticos para la maximización de la generación energética.

\section*{Situación problemática}
Los métodos tradicionales para determinar la inclinación y orientación de paneles fotovoltaicos se basan en fórmulas empíricas o aproximaciones estáticas, calculadas generalmente a partir de la latitud del sitio de instalación. Este enfoque asume condiciones de radiación solar constantes a lo largo del año y no contempla variaciones estacionales en la trayectoria solar, condiciones climáticas locales (nubosidad, temperatura ambiente, humedad) ni la presencia de obstáculos que generan sombras parciales sobre el arreglo fotovoltaico. Como consecuencia, las instalaciones dimensionadas con estos métodos suelen operar por debajo de su potencial energético real, generando pérdidas de eficiencia que no son detectadas ni corregidas por los procedimientos de cálculo convencionales. A esto se suma que, al tratarse de un problema con múltiples variables interrelacionadas y una función objetivo no lineal, los métodos analíticos clásicos presentan limitaciones para explorar de manera eficiente el espacio de soluciones posibles. Frente a esta situación, resulta necesario evaluar si técnicas de optimización basadas en algoritmos genéticos permiten superar dichas limitaciones y obtener configuraciones de inclinación y orientación que maximicen la energía captada en escenarios reales.
La inclinación (\textit{tilt}) y la orientación (azimut) de los paneles solares son variables determinantes en el rendimiento energético de una instalación. Una configuración inadecuada de estos parámetros puede provocar pérdidas significativas de generación, incluso cuando el resto del sistema —inversores, cableado, paneles— se encuentra correctamente dimensionado. Tradicionalmente, la determinación de estos ángulos se ha realizado mediante fórmulas empíricas o aproximaciones estáticas basadas en la latitud del lugar, sin considerar variaciones estacionales, condiciones climáticas locales, sombras del entorno o el comportamiento no lineal de la radiación solar a lo largo del año.
\section*{Problema}
Las instalaciones fotovoltaicas suelen configurarse con ángulos de inclinación y orientación fijos, determinados mediante criterios estáticos o generalizados que no consideran la variabilidad estacional de la radiación solar ni las condiciones específicas del entorno de instalación, lo que deriva en una captación energética subóptima.

\section*{Objetivos de la Investigación}
 
\subsection*{Objetivo General}
Desarrollar y evaluar un algoritmo genético capaz de optimizar la inclinación y orientación de paneles solares fotovoltaicos con el fin de maximizar la generación energética de una instalación determinada.
 
\subsection*{Objetivos Específicos}
\begin{enumerate}
  \item Analizar los fundamentos teóricos de la radiación solar, la geometría solar y su relación con la inclinación y orientación de los paneles fotovoltaicos.
  \item Identificar y sistematizar las variables ambientales y geográficas (latitud, estacionalidad, nubosidad, sombreado) que inciden en la generación energética de un sistema fotovoltaico.
  \item Diseñar un modelo de función objetivo que permita estimar la energía generada en función de los parámetros de inclinación y orientación de los paneles.
  \item Implementar un algoritmo genético que optimice los parámetros de inclinación y orientación a partir del modelo de función objetivo definido.
  \item Comparar los resultados obtenidos mediante el algoritmo genético con los valores calculados a través de métodos tradicionales o estáticos de configuración de paneles solares.
  \item Evaluar el desempeño del algoritmo genético en términos de precisión, tiempo de convergencia y mejora porcentual de la generación energética respecto de las configuraciones bases.
\end{enumerate}

\section*{Marco Teórico}
El rendimiento energético de un sistema fotovoltaico depende de la interacción entre variables astronómicas, geográficas, ambientales y propias del sistema de generación. La determinación de la inclinación y orientación óptimas requiere comprender previamente el comportamiento de la radiación solar y la geometría Sol-Tierra, ya que estas variables condicionan la cantidad de energía que puede captar un panel fotovoltaico. En esta sección se presentan los fundamentos teóricos que sustentan el desarrollo del modelo de optimización mediante algoritmos genéticos.

\subsection*{Sistema Fotovoltaico}
Un sistema fotovoltaico es un conjunto de dispositivos diseñados para convertir la energía proveniente de la radiación solar en energía eléctrica mediante el efecto fotovoltaico. Este proceso se produce cuando la luz solar incide sobre materiales semiconductores, generalmente silicio, generando el movimiento de electrones y, en consecuencia, una corriente eléctrica continua.
Los sistemas fotovoltaicos representan una de las principales tecnologías para el aprovechamiento de energías renovables debido a que permiten generar electricidad de forma limpia, silenciosa y sin emisiones directas de gases de efecto invernadero durante su operación.
Para garantizar el correcto funcionamiento de la instalación, intervienen diversos componentes que cumplen funciones específicas dentro del proceso de generación y suministro de energía eléctrica como lo son:
\begin{itemize}
    \item Paneles o módulos fotovoltaicos
    \item Inversor
    \item Estructura de soporte
    \item Elementos de protección y monitoreo
\end{itemize}
La cantidad de energía eléctrica producida por un sistema fotovoltaico depende de múltiples factores, entre ellos la intensidad de la radiación solar incidente, la temperatura de operación de los módulos, las condiciones atmosféricas, la presencia de sombras y, especialmente, la inclinación y la orientación de los paneles respecto de la trayectoria aparente del Sol. En este contexto, la correcta selección de dichos parámetros constituye un problema de optimización de gran relevancia, ya que una configuración inadecuada puede disminuir significativamente la energía captada durante la vida útil de la instalación.

\subsection*{Radiación Solar}
La radiación solar es la energía electromagnética emitida por el Sol que llega a la superficie terrestre tras atravesar la atmósfera. Constituye la fuente primaria de energía utilizada por los sistemas fotovoltaicos para generar electricidad. La cantidad de radiación disponible en un determinado lugar y momento depende de factores astronómicos, geográficos y atmosféricos, por lo que no permanece constante a lo largo del día ni del año.
Desde el punto de vista de los sistemas fotovoltaicos, la radiación solar puede clasificarse en tres componentes principales:
\begin{itemize}
    \item Radiación directa: corresponde a la energía que llega al panel siguiendo una trayectoria recta desde el Sol, sin sufrir desviaciones significativas en la atmósfera. Es el componente que aporta la mayor cantidad de energía cuando el cielo se encuentra despejado.
    \item Radiación difusa: es la radiación que ha sido dispersada por moléculas del aire, vapor de agua, aerosoles y nubes antes de alcanzar la superficie. Aunque posee menor intensidad que la radiación directa, continúa contribuyendo a la generación de energía, especialmente en días parcialmente nublados.
    \item Radiación reflejada (albedo): corresponde a la fracción de radiación que rebota sobre el suelo u otras superficies cercanas antes de llegar al panel fotovoltaico. Su contribución depende de las características del entorno, siendo mayor en superficies altamente reflectantes como nieve, arena clara o techos blancos.
\end{itemize}

\subsection*{Geometría Solar}
La geometría solar analiza la posición aparente del Sol con respecto a un punto determinado de la superficie terrestre. Esta posición varía continuamente debido a dos movimientos principales de la Tierra: la rotación sobre su eje, que provoca la sucesión del día y la noche, y la traslación alrededor del Sol, responsable de las estaciones del año.
Para describir la posición del Sol se emplean diversos ángulos astronómicos. Entre los más importantes se encuentran la declinación solar, el ángulo horario, la altura solar, el ángulo cenital y el azimut solar, los cuales constituyen la base para el cálculo del rendimiento energético de una instalación fotovoltaica.

\subsubsection*{Declinación Solar}
La declinación solar (d) es el ángulo formado entre los rayos del Sol y el plano del ecuador terrestre. Este parámetro varía continuamente durante el año como consecuencia de la inclinación del eje de rotación de la Tierra respecto de su plano orbital. Sus valores oscilan aproximadamente entre $+23,45^\circ$ durante el solsticio de junio y $-23,45^\circ$ durante el solsticio de diciembre. La declinación solar constituye uno de los parámetros fundamentales para describir la posición aparente del Sol y permite modelar la variación estacional de la radiación solar incidente. Se utiliza la fórmula de Cooper:
\[ d = 23,45^\circ \cdot \sin\left( \frac{360}{365} \cdot (n + 284) \right) \]
donde:
\begin{itemize}
    \item d = declinación solar en grados
    \item n = número del día del año
\end{itemize}
Esta fórmula define la variación de producción entre invierno y verano.

\subsubsection*{Ángulo Horario}
El ángulo horario ($\omega$) representa el desplazamiento angular del Sol respecto del meridiano local debido a la rotación de la Tierra. Se expresa en grados y equivale a $15^\circ$ por cada hora de diferencia respecto del mediodía solar. Su valor es negativo durante la mañana, igual a cero al mediodía solar y positivo durante la tarde.
\[ \omega = 15^\circ \cdot (HoraSolar - 12) \]
donde:
\begin{itemize}
    \item $\omega$ = Ángulo horario
    \item HoraSolar = Hora del día en formato 24 horas
\end{itemize}
Esta fórmula define la variación de producción a lo largo de un mismo día.

\subsubsection*{Altura Solar}
La altura solar ($\alpha$) es el ángulo formado entre los rayos solares y el plano horizontal. Este parámetro describe qué tan elevado se encuentra el Sol sobre el horizonte en un instante determinado y constituye una de las variables más importantes para estimar la radiación incidente sobre una superficie inclinada.
\[ \sin(\alpha) = \sin(L) \cdot \sin(d) + \cos(L) \cdot \cos(d) \cdot \cos(\omega) \]
donde:
\begin{itemize}
    \item $\alpha$ = Ángulo de altura solar
    \item L = Latitud del lugar
    \item d = Declinación solar
    \item $\omega$ = Ángulo horario
\end{itemize}
Esta fórmula determina la inclinación ideal del panel solar y el diseño de sombras.

\subsubsection*{Ángulo Cenital}
El ángulo cenital (Z) corresponde al ángulo formado entre la dirección del Sol y la vertical del lugar (cenit). Es complementario de la altura solar y se obtiene mediante la relación $Z = 90^\circ - \alpha$. Este parámetro resulta especialmente útil para estimar la cantidad de atmósfera que debe atravesar la radiación solar antes de alcanzar la superficie terrestre.
\[ Z = 90^\circ - \alpha \]
donde:
\begin{itemize}
    \item Z = Ángulo cenital
    \item $\alpha$ = Altura solar
\end{itemize}
Esta fórmula mide la pérdida por radiación difusa y el espesor de la atmósfera.

\subsubsection*{Azimut Solar}
El azimut solar (Az) describe la posición horizontal del Sol respecto de un punto de referencia geográfico. Este ángulo permite conocer la dirección desde la cual inciden los rayos solares y constituye un parámetro esencial para determinar la orientación óptima de un sistema fotovoltaico.
\[ \cos(Az) = \frac{\sin(\alpha) \cdot \sin(L) - \sin(d)}{\cos(\alpha) \cdot \cos(L)} \]
donde:
\begin{itemize}
    \item Az = Ángulo azimut solar
    \item $\alpha$ = Altura solar
    \item L = Latitud del lugar
    \item d = Declinación solar
\end{itemize}
Esta fórmula determina la orientación cardinal (hacia dónde debe apuntar el panel).

\end{document}